% 第5章图：两种消元次序(1,4,2,3)与(2,1,3,4)下原始图的变化与新作用域。
\begin{tikzpicture}[x=1mm,y=1mm,>=Stealth,
  v/.style={circle,draw=FigureStroke,fill=FigureBlueFill,line width=.7pt,minimum size=5.2mm,inner sep=0pt},
  cur/.style={v,fill=FigureYellowFill,draw=FigureYellow},
  e/.style={FigureStroke,line width=.7pt},
  fill edge/.style={FigureGreen,line width=.9pt,dashed},
  sc/.style={font=\figurefont\footnotesize,text=FigureInk},
  note/.style={font=\FigureNoteFont\fontsize{7.5}{9}\selectfont,text=FigureMuted}]
\def\P{27}

% ---------- row 1: order (1,4,2,3) ----------
\begin{scope}[yshift=36mm]
 \node[note,anchor=west] at (-4,20) {次序$(1,4,2,3)$};
 \begin{scope}[xshift=0mm]
  \node[cur] (a) at (0,6) {$1$}; \node[v] (b) at (8,12) {$2$}; \node[v] (c) at (8,0) {$3$}; \node[v] (d) at (16,6) {$4$};
  \draw[e] (a)--(b) (a)--(c) (b)--(c) (b)--(d) (c)--(d);
  \node[sc] at (8,-7) {$\{2,3\}$};
 \end{scope}
 \begin{scope}[xshift=\P mm]
  \node[v] (b) at (8,12) {$2$}; \node[v] (c) at (8,0) {$3$}; \node[cur] (d) at (16,6) {$4$};
  \draw[e] (b)--(c) (b)--(d) (c)--(d);
  \node[sc] at (8,-7) {$\{2,3\}$};
 \end{scope}
 \begin{scope}[xshift=2*\P mm]
  \node[cur] (b) at (8,12) {$2$}; \node[v] (c) at (8,0) {$3$};
  \draw[e] (b)--(c);
  \node[sc] at (8,-7) {$\{3\}$};
 \end{scope}
 \begin{scope}[xshift=3*\P mm]
  \node[cur] (c) at (8,0) {$3$};
  \node[sc] at (8,-7) {$\varnothing$};
 \end{scope}
\end{scope}
% ---------- row 2: order (2,1,3,4) ----------
\begin{scope}
 \node[note,anchor=west] at (-4,20) {次序$(2,1,3,4)$};
 \begin{scope}[xshift=0mm]
  \node[v] (a) at (0,6) {$1$}; \node[cur] (b) at (8,12) {$2$}; \node[v] (c) at (8,0) {$3$}; \node[v] (d) at (16,6) {$4$};
  \draw[e] (a)--(b) (a)--(c) (b)--(c) (b)--(d) (c)--(d);
  \node[sc] at (8,-7) {$\{1,3,4\}$};
 \end{scope}
 \begin{scope}[xshift=\P mm]
  \node[cur] (a) at (0,6) {$1$}; \node[v] (c) at (8,0) {$3$}; \node[v] (d) at (16,6) {$4$};
  \draw[e] (a)--(c) (c)--(d);
  \draw[fill edge] (a)--(d);
  \node[sc] at (8,-7) {$\{3,4\}$};
 \end{scope}
 \begin{scope}[xshift=2*\P mm]
  \node[cur] (c) at (8,0) {$3$}; \node[v] (d) at (16,6) {$4$};
  \draw[e] (c)--(d);
  \node[sc] at (8,-7) {$\{4\}$};
 \end{scope}
 \begin{scope}[xshift=3*\P mm]
  \node[cur] (d) at (16,6) {$4$};
  \node[sc] at (8,-7) {$\varnothing$};
 \end{scope}
\end{scope}
\foreach \y in {36,0}{\foreach \i in {0,1,2}{\draw[->,FigureMuted,line width=.6pt] ({\i*27+19},{\y+6})--({\i*27+23.6},{\y+6});}}
\end{tikzpicture}
